\documentclass[10pt]{article}

\usepackage[utf8]{inputenc}

\usepackage[T1]{fontenc}

\usepackage{amsmath}

\usepackage{amsfonts}

\usepackage{amssymb}

\usepackage[version=4]{mhchem}

\usepackage{stmaryrd}

\usepackage{graphicx}

\usepackage[export]{adjustbox}

\graphicspath{ {./images/} }



\begin{document}

Ф. с. является безусловным базисом во всех пространствах $L_{p}[0.1](1<p<\infty)$ и, более того, во всех рефлексивных пространствах Орлича (см. [5]). Если функция $f(t)$ принадлежит пространству $L_{p}[0,1]$, $1<p<\infty$, то имеют место неравенства



\[

A_{p}\|f\|_{p} \leqslant\left\|\left(\sum_{k=1}^{\infty} a_{k}^{2}(f) f_{k}^{2}(t)\right)^{1 / 2}\right\|_{p} \leqslant B_{p}\|f\|_{p},

\]



где $\|\cdot\|_{p}$ обозначает норму в пространстве $L_{p}[0,1]$, а постоянные $A_{p}>0$ и $B_{p}>0$ зависят лишь от $p$.



Ф. с. нашла важные приложения в различных вопросах анализа. В частности, с помощью әтой системы были построены базисы в пространствах $C^{1}\left(I^{2}\right)$ (см. [4]) и $A(D)$ (см. [5]). Здесь $C^{1}\left(I^{2}\right)$ - пространство всех непрерывно диффференцируемых на квадрате $I^{2}=[0,1] \times$ $\times[0,1]$ функциї $f(x, y)$ с нормой



\[

\|f\|=\max |f(x, y)|+\max \left|\frac{\partial f}{\partial x}\right|+\max \left|\frac{\partial f}{\partial y}\right| \text {, }

\]



а $A(D)$ - пространство всех функций $f(z)$ с нормой



\[

\|f\|=\max _{|z| \leqslant 1}|f(z)|

\]



аналитических в открытом круге $D=\{z:|z|<1\}$ комплексной плоскости и не и рерывных в замкнутом круге $\bar{D}=\{z:|z| \leqslant 1\}$. Вопросы о существовании базисов в пространствах $C^{1}\left(I^{2}\right)$ II $A(D)$ быліг поставлены C. Банахом [6].



JIum.: [1] F r a n k l i n P., "Math. Ann.", 1928, Bd 100,



\includegraphics[max width=\textwidth, center]{2023_07_09_9b2b68a8d99366363346g-1}

s k i Z., "Studia math.", 1963 , v. 23, N. 2 , p. $141-57$; [4] e r o



\includegraphics[max width=\textwidth, center]{2023_07_09_9b2b68a8d99366363346g-1(1)}

"Матем. сб.", 1974, т. $95, N_{2} 1$, c. 3-18; [6] B a n a c h S., Théorie des opérationes linéaires, Warsz., 1932.



ФРАТТИНИ ПОДГРӰППА - характеристическая подгруппа Ф $(G)$ группы $G$, определяемая как пересечение всех жаксимальных подгрупn $G$, если такие существуют; если же максимальных подгрупп в группе $G$ нет, то $G$ сама наз. своей Ф. п. Введена Дж. Фраттини [1]. Ф. п. состоит из тех и только тех элементов группы $G$. каждый из к-рых может быть удален из любой, содержащей его системы образующих групшы, T. e.



\[

\Phi(G)=\{x \mid x \in G,\{x, M\}=G \Rightarrow\{M\}=G\} .

\]



Конечная групша тогда и только тогда нильпотентна, когда ее коммутант содержится в ее Ф. п. Для любой конечной и любой полициклич. группы $G$ подгруппа $\Phi(G)$ нильпотентна.



Jum.: [1] F r a t t i n i G.. "Atti Accad. Lincei, Rend. (IV)",



\includegraphics[max width=\textwidth, center]{2023_07_09_9b2b68a8d99366363346g-1(2)}

ФРЕДГӦЛЬМА АЛЬТЕРНАТИВА - альтернативное утверждение, вытекающее из Фредгольма теорем. В случае линейного интегрального уравнения Фредгольма 2-го рода



\[

\varphi(x)-\lambda \int_{a}^{b} K(x, s) \varphi(s) d s=f(x), x \in[a, b],

\]



Ф. а. утверждает: либо уравнение (1) и сопряженное с ним уравнение



\[

\psi(x)-\bar{\lambda} \int_{a}^{b} \overline{K(s, x)} \psi(s) d s=g(x), x \in[a, b],

\]



имеют единственные решения $\varphi, \psi$, каковы бы ни были известные функции $f, g$, либо соответствующие однородные уравнения



\[

\begin{aligned}

& \varphi(x)-\lambda \int_{a}^{b} K(x, s) \varphi(s) d s=0, \\

& \psi(x)-\bar{\lambda} \int_{a}^{b} \overline{K(s, x)} \psi(s) d s=0

\end{aligned}

\]



имеют ненулевые решения, причем число линейно независимых решений конечно и одинаково для обоих уравнений. Во втором случае для того чтобы уравнение (1) имело решение, необходимо и достаточно, чтобы



\[

\int_{a}^{b} f(x) \overline{\psi_{k}(x)} d x=0, k=1, \ldots, n,

\]



где $\psi_{1}, \ldots, \psi_{n}$ - полная система линейно независимых решений уравнения $\left(2^{\prime}\right)$. При этом общее решение уравнения (1) имеет вид



\[

\varphi(x)=\varphi_{*}(x)+\sum_{k=1}^{n} c_{k} \varphi_{k}(x)

\]



где $\varphi_{k}$-какое-нибудь решение уравнения (1), $\varphi_{1}, \ldots$, $\varphi_{n}$-полная система линейно независимых решений уравнения $\left(1^{\prime}\right), r_{k}$-произвольные постоянные. Сходные утверждения имеют место и для уравнения (2).



Пусть $T$ - непрерывный линейный оператор, отображающий банахово пространство $E$ в себя; $E^{*}, T^{*}-$ соответствующие сопряженные пространство и оператор. Рассматриваются уравнения:



\[

\begin{gathered}

T(x)=y, x, y \in E, \\

T^{*}(g)=f, g, f \in E^{*}, \\

T(x)=0, x \in E, \\

T^{*}(g)=0, g \in E^{*} .

\end{gathered}

\]



Справедливость Ф. а. для оператора $T$ означает слөдующее: 1) либо уравнения (3) и (4) разрешимы, каковы бы ни были их правые части, и тогда их решения единственны; 2) либо однородные уравнения $\left(3^{\prime}\right)$ и $\left(4^{\prime}\right)$ пмеют одинаковое конечное число линейно независимых решений $x_{1}, \ldots, x_{n}$ и $g_{1}, \ldots, g_{n}$ соответственно; в этом случае для разрешимости уравнения (3) соответственно уравнения (4), необходимо и достаточно, чтобы $g_{k}(y)=0, k=1,2, \ldots, n$, соответственно $f\left(x_{k}\right)=0$, $k=1,2, \ldots, n ;$ при этом общее решение уравнения (3) дается равенством



\[

x=x^{*}+\sum_{k=1}^{n} c_{k} x_{k}

\]



а общее решение уравнения (4) - равенством



\[

g=g^{*}+\sum_{k=1}^{n} c_{k} g_{k}

\]



где $x^{*}$ (соответственно $g^{*}$ ) - какое-нибудь решение уравнения (3) (уравнения (4)), а $c_{1}, \ldots, c_{n}-$ произвольные постоянные.



Каждое из следующих двух условий необходимо и достаточно, чтобы для оператора $T$ имела место Ф. а. 1) Оператор $T$ представим в форме



\[

T=W+V

\]



где $W$ - опөратор, имеющий двусторонний непрерывный обратный, а $V$ - вполне непрерывный оператор, 2) оператор $T$ представим в форме



\[

T=W_{1}+V_{1},

\]



где $W_{1}$ - оператор, имеющий двусторонний непрерывный обратный, а $V_{1}$ - конечномерный оператор.



JІип.: [1] С м и р н о в B. II., Курс высшей математики, т. 4, ч. 1, 6 изд., М., 1974, [2] В ІІ а д и м и р о в В. С., Уравнения математической физики, 4 изд., М., 1981; А3] К а н т о p o-

в и ч Ј. В., А к и л о в Г. П., Фунццональный анализ, в и ч Л. В., А А и л о в Г. П., Функциональный анализ,

2 изд., М., 1977. ФРЕДГОЛЬМА ОIЕРАТОР - линейный оператор, отображающий одно локально выпуклое пространство в другое и ядерный оператор относительно Maккu monoлогии в этих Пространствах.

ФРЕДГОЛЬМА ТЕОРЕМЬІ д л я И. Бакушинский.

и т р а л ьны х у р а в не ни й:



Т е о р м а 1. Однородное уравнение



\[

\varphi(x)-\lambda \int_{a}^{b} K(x, s) \varphi(s) d s=0

\]





\end{document}